% GENERATED FILE. Edit canonical lessons, then run npm run generate:books.
% canonical-source-hash: fnv1a64:8fedbdd968b33dc9
% canonical-lessons: LA-C64-collum, LA-C64-umerus, LA-C64-genu, LA-C64-cor, LA-C64-sanguis

\chapter{Neck, Shoulder, Knee, Heart}
\label{ch:la-neck-shoulder-knee-heart}
\hlchaptermodality{64}

\begin{chapteropening}
\textbf{By the end of this chapter:} \emph{I can name the neck, a shoulder, a knee, the heart and blood.}

Complete the last lesson of chapter 64: I can name the neck, a shoulder, a knee, the heart and blood.
\end{chapteropening}

\section[collum, collī]{collum, collī — the neck}

\label{lesson:LA-C64-collum}



\begin{quote}
Before the new one: say the Latin for the tongue; a language, then the Latin for the nose.
\end{quote}

\subsection*{You'll want to know: collum, collī}
\textbf{collum, collī} — \textquotedblleft{}the neck\textquotedblright{}. A \emph{collar} goes round it.

\begin{cousinweb}
\textbf{collum} is inherited, usually traced to Proto-Indo-European \emph{*kwolso-}, \textquotedblleft{}neck\textquotedblright{} (EtymDB 2.1), perhaps from a root \emph{*kwel-}, \textquotedblleft{}to turn\textquotedblright{}.
\end{cousinweb}

\begin{grammarlens}[title={What you've built}]
One more word for this chapter.
\end{grammarlens}

\subsection*{Your turn}
\begin{itemize}
\raggedright
  \item \emph{Say it:} \emph{collum}
  \item \emph{Say it:} \emph{collum}, once more
  \item \emph{Say it:} say \emph{collum}
  \item \emph{Recall:} say the Latin for the tongue; a language, then the Latin for the nose, then say \emph{collum} again
\end{itemize}

\begin{tcolorbox}[breakable,colback=teal!4,colframe=teal!35!black,title={Before you move on}]
What does collum, collī mean? (\textquotedblleft{}The neck\textquotedblright{}.) Say it once more.
\end{tcolorbox}
\section[umerus, umerī]{umerus, umerī — the shoulder}

\label{lesson:LA-C64-umerus}



\begin{quote}
Before the new one: say the Latin for the nose, then the Latin for the neck.
\end{quote}

\subsection*{You'll want to know: umerus, umerī}
\textbf{umerus, umerī} — \textquotedblleft{}the shoulder\textquotedblright{}.

\begin{cousinweb}
\textbf{umerus} is inherited, usually traced to Proto-Indo-European \emph{*h3emsos}, \textquotedblleft{}shoulder\textquotedblright{} (EtymDB 2.1); Greek \emph{ōmos}, \textquotedblleft{}shoulder\textquotedblright{}, is a cousin (Lewis \& Short). English \emph{humerus}, the upper-arm bone, is the same word with a later \emph{h}.
\end{cousinweb}

\begin{grammarlens}[title={What you've built}]
One more word for this chapter.
\end{grammarlens}

\subsection*{Your turn}
\begin{itemize}
\raggedright
  \item \emph{Say it:} \emph{umerus}
  \item \emph{Say it:} \emph{umerus}, once more
  \item \emph{Say it:} say \emph{umerus}
  \item \emph{Recall:} say the Latin for the nose, then the Latin for the neck, then say \emph{umerus} again
\end{itemize}

\begin{tcolorbox}[breakable,colback=teal!4,colframe=teal!35!black,title={Before you move on}]
What does umerus, umerī mean? (\textquotedblleft{}The shoulder\textquotedblright{}.) Say it once more.
\end{tcolorbox}
\section[genū, genūs]{genū, genūs — a knee}

\label{lesson:LA-C64-genu}



\begin{quote}
Before the new one: say the Latin for the neck, then the Latin for the shoulder.
\end{quote}

\subsection*{You'll want to know: genū, genūs}
\textbf{genū, genūs} — \textquotedblleft{}a knee\textquotedblright{}. \emph{Genuflect}: to bend the knee.

\begin{cousinweb}
\textbf{genū} is inherited, usually traced to Proto-Indo-European \emph{*gonu}, \textquotedblleft{}knee\textquotedblright{} (EtymDB 2.1). English \emph{knee} is a cousin.
\end{cousinweb}

\begin{grammarlens}[title={What you've built}]
One more word for this chapter.
\end{grammarlens}

\subsection*{Your turn}
\begin{itemize}
\raggedright
  \item \emph{Say it:} \emph{genū}
  \item \emph{Say it:} \emph{genū}, once more
  \item \emph{Say it:} say \emph{genū}
  \item \emph{Recall:} say the Latin for the neck, then the Latin for the shoulder, then say \emph{genū} again
\end{itemize}

\begin{tcolorbox}[breakable,colback=teal!4,colframe=teal!35!black,title={Before you move on}]
What does genū, genūs mean? (\textquotedblleft{}A knee\textquotedblright{}.) Say it once more.
\end{tcolorbox}
\section[cor, cordis]{cor, cordis — the heart}

\label{lesson:LA-C64-cor}



\begin{quote}
Before the new one: say the Latin for the shoulder, then the Latin for a knee.
\end{quote}

\subsection*{You'll want to know: cor, cordis}
\textbf{cor, cordis} — \textquotedblleft{}the heart\textquotedblright{}. English \emph{cordial}, from the genitive \textbf{cordis}.

\begin{cousinweb}
\textbf{cor} is inherited, usually traced to Proto-Indo-European \emph{*krd-}, \textquotedblleft{}heart\textquotedblright{} (EtymDB 2.1). Greek \emph{kardia}, \textquotedblleft{}heart\textquotedblright{}, and English \emph{heart} are cousins (Lewis \& Short).
\end{cousinweb}

\begin{grammarlens}[title={What you've built}]
One more word for this chapter.
\end{grammarlens}

\subsection*{Your turn}
\begin{itemize}
\raggedright
  \item \emph{Say it:} \emph{cor}
  \item \emph{Say it:} \emph{cor}, once more
  \item \emph{Say it:} say \emph{cor}
  \item \emph{Recall:} say the Latin for the shoulder, then the Latin for a knee, then say \emph{cor} again
\end{itemize}

\begin{tcolorbox}[breakable,colback=teal!4,colframe=teal!35!black,title={Before you move on}]
What does cor, cordis mean? (\textquotedblleft{}The heart\textquotedblright{}.) Say it once more.
\end{tcolorbox}
\section[sanguis, sanguinis]{sanguis, sanguinis — blood}

\label{lesson:LA-C64-sanguis}



\begin{quote}
Before the new one: say the Latin for a knee, then the Latin for the heart.
\end{quote}

\subsection*{You'll want to know: sanguis, sanguinis}
\textbf{sanguis, sanguinis} — \textquotedblleft{}blood\textquotedblright{}.

\begin{cousinweb}
\textbf{sanguis} has no agreed origin; Lewis \& Short call it doubtful. EtymDB 2.1 records a proposal that ties it to Proto-Indo-European \emph{*h1esh2r}, \textquotedblleft{}blood\textquotedblright{}, but not everyone accepts it. English \emph{sanguine} comes from it.
\end{cousinweb}

\begin{grammarlens}[title={What you've built}]
One more word for this chapter.
\end{grammarlens}

\subsection*{Your turn}
\begin{itemize}
\raggedright
  \item \emph{Say it:} \emph{sanguis}
  \item \emph{Say it:} \emph{sanguis}, once more
  \item \emph{Say it:} say \emph{sanguis}
  \item \emph{Recall:} say the Latin for a knee, then the Latin for the heart, then say \emph{sanguis} again
\end{itemize}

\begin{tcolorbox}[breakable,colback=teal!4,colframe=teal!35!black,title={Before you move on}]
What does sanguis, sanguinis mean? (\textquotedblleft{}Blood\textquotedblright{}.) Say it once more.
\end{tcolorbox}
